% ============================================================
%  tikz_fig_adas_lr_shift.tex
%  長期 AD-AS：総需要シフトの効果
%
%  シナリオ: 拡張的財政・金融政策 → AD が右上にシフト
%  長期均衡: 実質 GDP は Y* に戻る（物価のみ P* → P** に上昇）
%  → 長期の「金融中立性」「古典派の二分法」の図解に使用
%
%  曲線の仕様:
%    AD  : P = 10/Y + 1.5  （シフト前）
%    AD' : P = 14/Y + 1.5  （シフト後）
%    LRAS: Y = Y* = 5.0
%    均衡1: (5, 3.5)  ←シフト前
%    均衡2: (5, 4.3)  ←シフト後
% ============================================================
\documentclass[tikz, dvipdfmx]{standalone}
\usepackage{tikz}
\usetikzlibrary{arrows.meta}
\definecolor{gdpblue} {RGB}{30,  90,  180}
\definecolor{trendred}{RGB}{190, 40,  40}

\begin{document}
\begin{tikzpicture}[
  >=Stealth,
  thick,
]

% ── 軸 ──────────────────────────────────────────────────────
\draw[->] (0,0) -- (9.8,0)
  node[right, font=\small] {$Y$ (Output)};
\draw[->] (0,0) -- (0,7.8)
  node[above, font=\small] {$P$ (Price level)};

% ── LRAS（垂直線） ────────────────────────────────────────
\draw[trendred, very thick] (5,0) -- (5,7.2)
  node[above, font=\small, trendred] {LRAS};

% ── AD（シフト前, 実線） ─────────────────────────────────
\draw[gdpblue, very thick]
  plot[domain=1.72:9.3, samples=120, smooth] (\x, {10/\x + 1.5})
  node[right, font=\small, gdpblue] {AD};

% ── AD'（シフト後, 破線） ────────────────────────────────
%    P = 14/Y + 1.5; 新均衡 (5, 4.3)
\draw[gdpblue, very thick, dashed]
  plot[domain=2.2:9.3, samples=120, smooth] (\x, {14/\x + 1.5})
  node[right, font=\small, gdpblue] {AD$'$};

% ── シフト矢印（ADからAD'へ） ────────────────────────────
%    曲線の中間あたりに矢印
\draw[->, thick, gdpblue!70]
  (3.3, {10/3.3 + 1.5 + 0.2}) -- (3.3, {14/3.3 + 1.5 - 0.2});

% ── 均衡点1（シフト前） ──────────────────────────────────
\fill[black]  (5, 3.5) circle[radius=2.5pt];
\node[left, font=\scriptsize, xshift=-2pt] at (5,3.5) {$E$};

% ── 均衡点2（シフト後） ──────────────────────────────────
\fill[black]  (5, 4.3) circle[radius=2.5pt];
\node[left, font=\scriptsize, xshift=-2pt] at (5,4.3) {$E'$};

% ── P* への水平点線（シフト前） ──────────────────────────
\draw[dashed, gray!70, thin]
  (0, 3.5) node[left, font=\small, black] {$P^*$} -- (5, 3.5);

% ── P** への水平点線（シフト後） ─────────────────────────
\draw[dashed, gray!70, thin]
  (0, 4.3) node[left, font=\small, black] {$P^{**}$} -- (5, 4.3);

% ── 物価上昇の縦矢印 ─────────────────────────────────────
\draw[<->, semithick, black!60]
  (0.6, 3.5) -- (0.6, 4.3);

% ── Y* ───────────────────────────────────────────────────
\node[below, font=\small] at (5,0) {$Y^*$};

% ── 「実質 GDP は不変」の注釈 ────────────────────────────
\node[
  align  = center,
  font   = \scriptsize,
  black!70,
] at (7.5, 1.5) {$\Delta Y = 0$ in LR\\(only $P$ rises)};

\end{tikzpicture}
\end{document}
